\documentclass[10pt]{article}
\usepackage[T1]{fontenc}
\usepackage[utf8]{inputenc}
\usepackage[french]{babel}
\usepackage{amsmath,amssymb,mathrsfs,mathpazo}
\usepackage{geometry}
\geometry{a4paper,margin=2cm}
\usepackage{graphicx,array,tabularx,tabularray,colortbl,multirow,multicol,booktabs,enumitem,pifont,ulem,fancybox,diagbox,cancel}
\usepackage{pstricks,pstricks-add,pst-plot,pst-tree,pst-node,pst-text,pst-func,pst-math,pst-bspline,pst-3dplot,pst-eucl,pst-solides3d}
\let\pscalculate\undefined
\usepackage{pst-bezier}
\usepackage{tikz,pgfplots,tkz-tab}
\usetikzlibrary{arrows,arrows.meta,patterns,calc,positioning,trees,decorations.markings,shapes.geometric,shapes.misc,shapes.arrows,intersections,decorations.pathmorphing,decorations.pathreplacing,backgrounds,shadows,babel}
\pgfplotsset{compat=1.18}
\usepgfplotslibrary{fillbetween}
\usepackage[np]{numprint}
\usepackage{listings,makecell,pas-tableur,twoopt}
\let\barre\undefined
\usepackage[pstricks]{tablvar}
\definecolor{darkgreen}{RGB}{0,100,0}
\providecommand{\red}{\color{red}}
\providecommand{\blue}{\color{blue}}
\providecommand{\green}{\color{green}}
\providecommand{\black}{\color{black}}
\providecommand{\white}{\color{white}}
\providecommand{\gray}{\color{gray}}
\providecommand{\cyan}{\color{cyan}}
\providecommand{\magenta}{\color{magenta}}
\providecommand{\yellow}{\color{yellow}}
\providecommand{\violet}{\color{violet}}
\providecommand{\orange}{\color{orange}}
\providecommand{\euro}{\texteuro}
\providecommand{\EUR}{\texteuro}
\providecommand{\ieme}{\textsuperscript{e}}
\providecommand{\degres}{\ensuremath{^\circ}}
\providecommand{\degree}{\ensuremath{^\circ}}
\providecommand{\vv}[1]{\overrightarrow{#1}}
\providecommand{\Coord}[1]{#1}
\providecommand{\decosix}{}
\providecommand{\starredbullet}{\ensuremath{\bullet}}
\pagestyle{empty}
\setlength{\parindent}{0pt}
\usepackage{adjustbox}\DeclareUnicodeCharacter{2236}{:}\DeclareUnicodeCharacter{2190}{\ensuremath{\leftarrow}}
\providecommand{\R}{\mathbb{R}}
\providecommand{\N}{\mathbb{N}}
\providecommand{\Z}{\mathbb{Z}}
\providecommand{\Q}{\mathbb{Q}}
\providecommand{\D}{\mathbb{D}}
\providecommand{\C}{\mathbb{C}}
\providecommand{\e}{{\mathrm{e}}}
\providecommand{\vect}[1]{\overrightarrow{#1}}
\providecommand{\vectt}[1]{\overrightarrow{\mathrm{#1}}}
\providecommand{\vvt}[1]{\overrightarrow{\mathrm{#1}}}
\providecommand{\np}[2][]{{#2~#1}}
\providecommand{\vv}[1]{{\overrightarrow{#1}}}
\renewcommand{\d}{{\mathrm{d}}}
\providecommand{\degres}{{^\circ}}
\providecommand{\degre}{{^\circ}}
\providecommand{\blue}{\color{blue}}
\providecommand{\red}{\color{red}}
\providecommand{\enskip}{\quad}
\providecommand{\llbracket}{\lbrack\!\lbrack}
\providecommand{\rrbracket}{\rbrack\!\rbrack}
\providecommand{\euro}{{\text{€}}}
\providecommand{\no}{{\text{n°}}}
\providecommand{\interval}[2]{{[#1;#2]}}
\providecommand{\intervFO}[2]{{[#1;#2[}}
\providecommand{\dotfill}{\ldots}

\providecommand{\ds}{\displaystyle}
\providecommand{\textding}[1]{#1}
\providecommand{\Oij}{\(\mathrm{O};\vec{\imath},\vec{\jmath}\)}
\providecommand{\Oijk}{\(\mathrm{O};\vec{\imath},\vec{\jmath},\vec{k}\)}
\providecommand{\Ouv}{\(\mathrm{O};\vec{u},\vec{v}\)}
\providecommand{\intervFF}[2]{\ensuremath{[#1;#2]}}
\renewcommand{\euro}{\ifmmode\text{\texteuro}\else\texteuro\fi}
\DeclareUnicodeCharacter{00B0}{\ifmmode^{\circ}\else\textdegree\fi}
\DeclareUnicodeCharacter{20AC}{\ifmmode\text{\texteuro}\else\texteuro\fi}
% Énoncés : Éducation nationale. Transcription et corrigés : APMEP, Denis Vergès et contributeurs. Découpage et classement : CoopMaths, Sébastien Lozano, Éric Elter et contributeurs. Adaptation de la présentation : Maths Explorer.
% CC BY-SA : https://coopmaths.fr/www/about/
% https://coopmaths.fr/alea/static/e3c/2020/tex/e3c_2020_00_sujet7_3_cor.tex
\begin{document}
\noindent\textit{APMEP / CoopMaths. CC BY-SA. Présentation adaptée par Maths Explorer.}\par\medskip
\subsection*{1. Calculer les termes d'indice 3 des suites $(u_n)$ et $(v_n)$.}
	
	\[
	\begin{array}{|c|c|c|}
		\hline
			n & u_n & v_n \\
		\hline
		0 & -4 & 0 \\
			\hline
		1 & 2 & 3,5 \\
			\hline
		2 & 4 & 5,25 \\
			\hline
		3 & 5 & 6,125 \\
			\hline
	\end{array}
	\]
	
	\subsection*{2. On s'intéresse aux variations de la suite $(u_n)$. Pour cela, on considère la fonction $f$ définie sur $[0; +\infty[$ par : $f(x) = \dfrac{8x - 4}{x + 1}$}
	
	\subsubsection*{a. Démontrer que la fonction $f$ est croissante sur $[0; +\infty[$.}
	
	Sur $[0; +\infty[$, la fonction est dérivable et sur cet intervalle :
	\[
	f'(x) = \dfrac{8(x + 1) - (8x - 4)}{(x + 1)^2} = \dfrac{8x + 8 - 8x + 4}{(x + 1)^2} = \dfrac{12}{(x + 1)^2}
	\]
	
	$f'(x) > 0$ car quotient de deux nombres supérieurs à zéro ; la fonction $f$ est donc strictement croissante sur l'intervalle $[0; +\infty[$.
	
	\subsubsection*{b. En déduire la monotonie de la suite $(u_n)$.}
	
	De la croissance de la fonction résulte la stricte croissance de la suite $(u_n)$.
	
	\subsection*{3. On considère l'affirmation suivante : « pour tout entier $n$, $u_n < v_n$ ». Camille pense que cette affirmation est vraie alors que Dominique pense le contraire. Pour les départager, on réalise le programme suivant écrit en langage Python :}
	
\begin{center}
\begin{tabularx}{0.5\linewidth}{|X|}\hline
\quad \texttt{def algo( )  :}\\
\quad \texttt{n = 0}\\
\quad \texttt{u = -4}\\
\quad \texttt{v = 0}\\
\quad \texttt{while u < v}\\
\qquad \texttt{n = n+1}\\
\qquad \texttt{u = (8*n - 4)/(n + 1)}\\ 
\qquad \texttt{v = 0,5*v + 3,5}\\
\quad \texttt{return(n)}\\ \hline
\end{tabularx}
\end{center}

	
	Le programme dit que $u_{11} > v_{11}$, donc Dominique a raison.\\ Effectivement $u_{11} = \dfrac{84}{12} = 7$ et $v_{11} = 6,99658203125$.

\end{document}
